\subsection{Actual finite images at the coefficient prime eleven}

The following finite quotients refine the weight-nine certificate.
They are quotients of the actual groups, with all source-lifting
and arithmetic lower bounds included. For a positive integer $c<11$,
let $L_{\le c}$ be the free nilpotent Lie algebra on $X,Y$ over
$\mathbb F_{11}$, truncated above degree $c$, and put
\[
 Q_c=\exp(L_{\le c}).                                      \tag{B.F1}
\]
Here multiplication is the Campbell--Hausdorff polynomial. Every
one of its denominators through class ten is prime to eleven.
Thus $Q_c$ is the finite free exponent-eleven group of class $c$
on $x=\exp X,y=\exp Y$. Equivalently it is the quotient of the
free pro-eleven group $P$ by the closed subgroup generated by
all eleventh powers and $\gamma_{c+1}P$. The two descriptions
are identified by these named generators. The logarithm and
exponential polynomials are mutually inverse; their primitive
terms are precisely the free Lie terms below degree eleven.
This proves the stated identification, including its integral
coefficient reduction. Since the quotient is characteristic,
the normalized genuine action of $\GTz$ on $\widehat F(x,y)$
gives a continuous homomorphism
\[
 \rho_c:\GTz\longrightarrow\operatorname{Aut}(Q_c).
                                                               \tag{B.F2}
\]
It is obtained by restriction from the action on a finite
characteristic quotient of $\widehat B$: the pure subgroup
$F$ is characteristic, and its power and lower-central kernels
are characteristic in $F$.

\begin{result}{B.2.2}{Proved}{Exact finite images through class ten at eleven}
For $3\le c\le10$ we have
\[
 \rho_c(G_{\mathbb Q})=\rho_c(\GT),\qquad
 \rho_c(\ker\Delta)=\rho_c(\GTz).
                                                               \tag{B.F3}
\]
Their orders are given below. The second and third columns are
the orders of the multiplier-one images; every full image has
order ten times the corresponding entry.
\[
\begin{array}{c|r|r|r|r}
 c&|\rho_c(\mathcal H)|&|\rho_c(\GT\cap\mathcal H)|
   &|\rho_c(\GTz)|&|\rho_c(\GT)|\\\hline
 3&1&1&10&10\\
 4&11&11&110&110\\
 5&11&11&110&110\\
 6&121&121&1210&1210\\
 7&121&121&1210&1210\\
 8&14641&1331&146410&13310\\
 9&161051&14641&1610510&146410\\
 10&2357947691&161051&23579476910&1610510
\end{array}                                                   \tag{B.F4}
\]
Thus class eight is the first separating quotient in this
specified finite sequence. These quotients give no separator
between $G_{\mathbb Q}$ and $\GT$.
\end{result}
\begin{proof}
We first describe the exact filtration and its upper bounds.
For a normalized multiplier-one automorphism, its pair coordinate
$f$ is recovered, modulo degree $c$, from
$y\mapsto f^{-1}yf$ in $Q_c$. Indeed, the leading Lie centralizer
of $Y$ consists of its scalar multiples. A conjugator ambiguity
below degree $c$ is therefore a power of $y$; its exponent is
zero because $f$ has zero abelianization. The central degree-$c$
term acts trivially and is not included in this assertion.
Consequently the kernel of passage from class $n+1$ to class $n$
is identified with a subspace of the degree-$n$ free Lie algebra,
by the first term of $f$, for $2\le n\le9$. Composition in this
kernel is addition, and cyclotomic conjugation by multiplier
$a$ acts on that subquotient as $a^n$.

Write $Z=-X-Y$. A possible first term $h$ has the following
four exact linear constraints:
\[
\begin{aligned}
 h(X,Y)+h(Y,X)&=0,\\
 [X,h(X,Y)]-[X+Y,h(Z,Y)]&=0,\\
 [h(X,Y),Y]-[-h(Z,X),X+Y]&=0,\\
 h(X,Y)+h(Y,Z)+h(Z,X)&=0.
\end{aligned}\tag{B.F5}
\]
The middle two equations express the special derivation and
its cyclic conjugate; the first and fourth express its outer
invariance under the geometric $\mathfrak S_3$. They follow by
taking the first term of the actual peripheral and symmetry
identities. Conversely they say that the special outer
class of $D_h(X)=0,D_h(Y)=[h,Y]$ is fixed by that same
$\mathfrak S_3$, with the inner adjustments explicitly
computed in the proof of B.2.1. Thus (B.F5) describes the
entire leading special invariant space, not merely a selected
family of elements of it.

For the pentagon we use the actual lower-central comparison
for $P\Gamma_{0,5}=P_4/Z(P_4)$. In degrees under consideration
its Lie algebra has the presentation
\[
 \mathfrak t_4/\langle\textstyle\sum_{i<j}t_{ij}\rangle
       =L(a,b,c)\rtimes L(u,v),                            \tag{B.F6}
\]
where $a=t_{14},b=t_{24},c=t_{34},u=t_{12},v=t_{23}$ and
$t_{13}=-u-v-a-b-c$. The two derivations on the free ideal are
\[
\begin{array}{lll}
 D_u(a)=[a,b],&D_u(b)=[b,a],&D_u(c)=0,\\
 D_v(a)=0,&D_v(b)=[b,c],&D_v(c)=[c,b].
\end{array}                                                \tag{B.F7}
\]
These formulas are the degree-two terms of the pure-braid
conjugation maps; the group presentation is the same actual
forgetting split extension used in B.5. To check that no
rational-completion substitution has occurred, a split extension
$F_3\rtimes F_2$ with this trivial action on $F_3^{\rm ab}$
has
$\gamma_n(F_3\rtimes F_2)=\gamma_nF_3\rtimes\gamma_nF_2$.
The commutator inclusion
$[\gamma_iF_2,\gamma_jF_3]\subseteq\gamma_{i+j}F_3$
follows by induction from the displayed action and the
three-subgroup commutator identity. This gives one containment;
the two factors' lower-central groups give the reverse one.
Hence every graded piece is the sum of the two free integral
Lie pieces, and is torsion free. Collecting commutators below
class eleven, or applying the same Campbell--Hausdorff
polynomials over $\mathbb Z_{11}$ and then reducing, retains
these pieces modulo eleven. This proves the exact finite
lower-central comparison needed in (B.F6).

The first term of the actual pentagon is therefore
\[
\begin{aligned}
 \mathcal P_n(h)={}&h(t_{12},t_{23}+t_{24})
       +h(t_{13}+t_{23},t_{34})-h(t_{23},t_{34})\\
       &-h(t_{12}+t_{13},t_{24}+t_{34})-h(t_{12},t_{23}).
\end{aligned}                                               \tag{B.F8}
\]
A genuine pentagon element maps to its kernel. Lower pair
coefficients that are multiples of eleven disappear under the
actual coefficient reduction, so the same assertion applies
to the finite filtration kernels, not only to elements with
integrally zero lower terms.

The program \texttt{check\_low\_weight\_images.py} expands every
standard Lyndon bracket in the free associative algebra. It
forms all the coefficient equations in (B.F5), reduces the
full matrix over $\mathbb F_{11}$, and evaluates (B.F8) by
(B.F6)--(B.F7). In the latter computation the free-$L(u,v)$
terms cancel literally; the remaining free-$L(a,b,c)$
polynomials are compared coefficient by coefficient. The
result is
\[
\begin{array}{c|rrrrrrrr}
 n&2&3&4&5&6&7&8&9\\\hline
 \dim L_n(X,Y)&1&2&3&6&9&18&30&56\\
 \dim V_n\text{ from (B.F5)}&0&1&0&1&0&2&1&4\\
 \dim\ker(\mathcal P_n|V_n)&0&1&0&1&0&1&1&1.
\end{array}                                                 \tag{B.F9}
\]
Every matrix, basis vector and resulting kernel is retained
in its JSON certificate. In each degree we also compute an
exact rational nullspace using rows independent modulo eleven.
After clearing denominators prime to eleven, every such vector
is checked against every integer row. Their reductions have
the full dimension in (B.F9). Thus the integral leading
lattice surjects onto $V_n$; this is an executed integral
lifting check, not an inference from dimensions in different
characteristics.

Every integral vector just obtained is the leading term of
an actual full gentle element. The construction in B.2.1
works for each $n\ge2$: choose lifts of the two peripheral
conjugators, cancel each higher product error using the
surjection
\[
 (u,v)\longmapsto[u,Y]+[v,Z]\quad\hbox{onto }L_m(X,Y),
\]
then lift the invariant outer class through the pro-eleven
special-automorphism kernel by coprime complement conjugacy.
The same cusp normalization and the same refereed relative
framed theorem R1 produce a full profinite element in
$\mathcal H$. These operations retain the first term
modulo eleven. This proves the lower bound $\dim V_n$ in
all the gentle filtration kernels.

We now give actual arithmetic elements supplying every
pentagon lower bound. For $m=3,5,7,9$ put
\[
 u_m=\prod_{a=1}^{10}(1-\zeta_{11}^a)^{e_{a,m}},\qquad
 e_{a,m}\in\{0,\ldots,10\},\quad
 e_{a,m}\equiv a^{m-1}\pmod{11}.                  \tag{B.F10}
\]
This representative has the same eleventh-power class as the
product with the unreduced exponents $a^{m-1}$.
Modulo eleventh powers, changing $\zeta_{11}$ to
$\zeta_{11}^b$ multiplies the class of $u_m$ by $b^{1-m}$.
Its Kummer extension therefore has a $C_{11}$ kernel on
which the cyclotomic quotient acts by $b^m$. The following
literal finite-field checks show that this kernel is
nontrivial:
\[
\begin{array}{c|r|r|r|r}
 m&q&\zeta_{11}\bmod q&u_m\bmod q&u_m^{(q-1)/11}\bmod q\\\hline
 3&23&2&8&18\\
 5&23&2&9&12\\
 7&23&2&12&6\\
 9&67&64&26&15.
\end{array}                                                 \tag{B.F11}
\]
Each last entry differs from one. The kernel thus lies in
the derived subgroup of the Galois group of the Kummer
extension: choose $b$ with $b^m\ne1$, and its commutator
with a kernel generator generates that $C_{11}$. Lifting
these two elements to $G_{\mathbb Q}$ and taking their
commutator gives an element of $G_{\mathbb Q^{\rm ab}}$
whose $m$th Soul\'e character is a unit modulo eleven.
The character is precisely the one specified by (B.F10),
by \cite[equation (1.2)]{NSW2017} at level $n=1$.
Reindexing the product at every level $11^n$ also proves
its cyclotomic conjugation weight $m$.

We retain the comparison to the pair coordinate. In the
completed metabelian group of $P$, let $C=[x,y]$, and write
$f=C^{a_f}$ in the free rank-one module over
$\Lambda=\mathbb Z_{11}[[\bar x-1,\bar y-1]]$. The Fox
resolution proves this rank-one assertion: the kernel of
$\Lambda^2\to\Lambda$, $(a,b)\mapsto(\bar x-1)a+
(\bar y-1)b$, is generated by
$(1-\bar y,\bar x-1)$, since these two variables form a
regular sequence. For the normalized action
$x\mapsto x,y\mapsto f^{-1}yf$, direct multiplication gives
\[
 \alpha(C)=C^{B_\alpha},\qquad
 B_\alpha=1+(\bar x-1)(\bar y-1)a_f.                     \tag{B.F12}
\]
The multiplier-one specialization of the arithmetic beta
series is this action on the commutator module
\cite[Section 0.2]{Nakamura2019}. The explicit formula
\cite[equation (1.1)]{NSW2017}, with
$\bar x=e^U,\bar y=e^V$, is
\[
 \log B_\sigma=
 \sum_{m\ge3,\ m\ {\rm odd}}
 \frac{\chi_m(\sigma)}{11^{m-1}-1}
 \sum_{\substack{i+j=m\\i,j\ge1}}\frac{U^iV^j}{i!j!}.   \tag{B.F13}
\]
For the four weights used, its $U^{m-1}V$ coefficient is
$\chi_m/((11^{m-1}-1)(m-1)!)$, a unit multiple of
$\chi_m$. This proves the required integral normalization.
In particular $\chi_m$ modulo eleven depends only on the
actual normalized action through class $m+1$.

Here is an explicit projection that removes all smaller
weights without assuming a section of the cyclotomic
character in $G_{\mathbb Q}$. Work first over
$\mathbb Z_{11}$ through class ten, and let $H$ be the
compact image of $G_{\mathbb Q^{\rm ab}}$ in the normalized
IA automorphisms. This is a pro-eleven group of nilpotency
class less than eleven. Campbell--Hausdorff logarithms
identify it with a closed Lie subring of the filtered
IA Lie ring. The assertion about closed subgroups follows
by solving the finite Campbell--Hausdorff formulas for
addition and brackets, weight by weight: only denominators
$1,\ldots,10$ occur, their inverse powers are defined by
continuous $\mathbb Z_{11}$ powering, and closure retains
the resulting limits. Thus this Lie subring is closed
under $\mathbb Z_{11}$ scalars.

Choose an arithmetic element with eleven-adic multiplier
a primitive Teichm\"uller tenth root $\zeta$. Its image
has a procyclic closure, whose kernel over the multiplier
is pro-eleven. The prime-to-eleven component of that
procyclic image has order ten; taking the corresponding
profinite power of the arithmetic element realizes this
component. Denote its action on $H$ by $T$. We use the
order-ten action in the truncated image; we assert no
order-ten element of $G_{\mathbb Q}$. On the degree-$n$
quotient of $\log H$, $T$ acts as $\zeta^n$.
For each of the four weights apply
\[
       e_m=\frac1{10}\sum_{i=0}^9\zeta^{-mi}T^i
             \quad\hbox{to }\log\sigma.                 \tag{B.F14}
\]
All its coefficients lie in $\mathbb Z_{11}$, so its
exponential is still an actual arithmetic image. The
nonzero differences $\zeta^m-\zeta^n$, $2\le n<m$, are
units. Comparing successive filtration quotients proves
$e_m\log\sigma$ has no term below weight $m$. The
Soul\'e character is a continuous homomorphism of $H$,
hence a linear functional on its Lie ring, with weight
$m$. It is unchanged by $e_m$; (B.F11)--(B.F13) prove
that the resulting degree-$m$ pair term is nonzero
modulo eleven. We have therefore obtained genuine
arithmetic generators for the one-dimensional kernels
in weights $3,5,7,9$.

For weight eight use the commutator of the weight-three
and weight-five elements. With
$\partial_h(X)=0,\partial_h(Y)=[Y,h]$, its pair term is
\[
        \{h,k\}=[h,k]+\partial_h(k)-\partial_k(h).
                                                               \tag{B.F15}
\]
This formula follows by evaluating the commutator of the
two derivations on $Y$. The exact program computes it
on the bases of (B.F9) and obtains six times the displayed
weight-eight basis modulo eleven. Hence it is nonzero.
All the pentagon upper bounds are now attained by actual
arithmetic elements. Descending through the finite
filtration proves equality of the multiplier-one images
of $G_{\mathbb Q}$ and $\GT$, and gives their orders as
$11^{\sum_{n<c}\dim\ker\mathcal P_n}$. The analogous
gentle order is $11^{\sum_{n<c}\dim V_n}$.

For precision, the kernel of multiplier modulo eleven
in each full finite image is already the image of its
actual multiplier-one subgroup. To prove this, let $K_n$
be its degree-$n$ filtration. Commuting with the arithmetic
Teichm\"uller element multiplies $K_n/K_{n+1}$ by
$\zeta^n-1$, a unit for $2\le n\le9$. Every such
commutator has actual multiplier one. Induction gives
$K_n\subseteq\rho_c(\ker\chi)K_{n+1}$ and hence the
claim. The arithmetic multiplier image is all of
$\mathbb F_{11}^\times$, of order ten. This proves every
full order and the first equality of (B.F3).

Finally the R1 witness of A.2 has $\Delta=1$ and maps to
one under every $\rho_c$. Its pair coordinate lies in
$R_{\rm nil}(F'F^3)$; the map $F'F^3\to P$ is onto,
because cubing is invertible on each pro-eleven cyclic
generator. Restriction of every finite eleven-group
quotient of $P$ is a finite nilpotent quotient of that
subgroup, so the pair coordinate maps to one in $P$.
Multiplying an arbitrary full gentle lift by this witness
changes its parity and preserves its finite image. This
proves the second equality of (B.F3).
\end{proof}

\begin{result}{B.2.3}{Proved}{An integral weight-seven separating polynomial}
Let brackets indexed by a Lyndon word mean its standard
Lyndon bracketing, and put
\[
\begin{aligned}
 h_7={}&-3[xxxyxyy]-3[xxxyyxy]-[xxyxyxy]
       +3[xxyxyyy]\\
      &+3[xxyyxyy]+3[xxyyyxy]+[xyxyxyy].
\end{aligned}
\tag{B.F16}
\]
This is the first term of an actual full gentle
multiplier-one element. Its pentagon has nonzero image
\[
          \mathcal P_7(h_7)\longmapsto
                    160((34)-(12))\in\mathbb Z[S_4]
                                                               \tag{B.F17}
\]
under the infinitesimal representation used in B.2.1.
Its square has parity zero and nonzero pentagon, providing
an additional actual witness for $\GT\subsetneq\ker\Delta$.
\end{result}
\begin{proof}
The integer source-nullspace check verifies all four
identities (B.F5) for (B.F16), coefficient by coefficient.
The full peripheral, coprime and R1 lifts in the preceding
proof therefore apply to this precise integral polynomial.
Substitution in (B.F8) under $t_{ij}\mapsto(ij)-C/6$
produces (B.F17), by exact integer arithmetic. The scalar
$160$ is six modulo eleven, so the defect survives the
actual finite quotient. Squaring adds the two first terms,
retains a nonzero coefficient modulo eleven, and kills
$\Delta$ because it is a homomorphism to $\mathbb F_2$.
\end{proof}

\begin{result}{B.2.4}{Proved}{Odd-prime pronilpotent actions detect no quadratic arithmetic annihilator}
For every odd prime $\ell$, let $P_\ell$ be the free
pro-$\ell$ quotient of $\widehat F(x,y)$ and let
\[
 r_{\rm odd}:\GTz\longrightarrow
       \prod_{\ell\ {
m odd}}\operatorname{Aut}(P_\ell)
                                                               \tag{B.F18}
\]
be the continuous homomorphism given by the genuine normalized
actions. Every continuous quadratic character of $\GTz$
which factors through (B.F18) and vanishes on $G_{\mathbb Q}$
is zero. In particular the same assertion holds for each
finite characteristic odd-prime group quotient of the free
group, and every member of $\mathscr X$ detected by
(B.F18) is zero.
\end{result}
\begin{proof}
Put $S=r_{\rm odd}(\GTz)$, a compact subgroup of the
displayed product. The action on
$P_\ell/\Phi(P_\ell)=\mathbb F_\ell^2$ is multiplication
by the reduction of the multiplier. Thus it gives an exact
sequence of compact groups
\[
 1\longrightarrow K\longrightarrow S
     \xrightarrow{\,a\,}
       \prod_{\ell\ {
m odd}}\mathbb F_\ell^\times
     \longrightarrow1.                                    \tag{B.F19}
\]
Surjectivity follows from the arithmetic subgroup: the Galois
action on the roots of unity gives the full cyclotomic group
$\widehat{\mathbb Z}^{\times}$, and reduction at the odd
primes maps it onto the displayed product by the Chinese
remainder theorem and compactness.

The kernel of
$\operatorname{Aut}(P_\ell)\to
\operatorname{GL}_2(\mathbb F_\ell)$ is pro-$\ell$.
We recall the finite argument that gives exactly this
property. On a finite $\ell$-group generated by two
elements, an automorphism acting identically on its
Frattini quotient acts identically on every graded term
of the lower exponent-$\ell$ central series: those
terms are generated by the original degree-one classes
under commutators and $\ell$th powers. Successive
kernels of the action on the series quotients embed,
by the difference on the two generators, in additive
$\mathbb F_\ell$-vector spaces. The last kernel is
trivial because the series terminates. The Frattini
kernel therefore has $\ell$-power order. Applying this
to the cofinal characteristic finite quotients of
$P_\ell$ proves the assertion for its automorphism
group.

Consequently $K$ is a closed subgroup of a product of
pro-$\ell$ groups with $\ell$ odd. Every finite
continuous quotient of $K$ has odd order. Indeed, its
open kernel contains the intersection with an open
normal subgroup of the product specified at finitely
many primes; the corresponding finite product and
its subgroup image both have odd order. A character
$S\to\mathbb F_2$ therefore kills $K$ and factors
continuously through $a$. The arithmetic image
surjects onto the target of $a$, so arithmetic
vanishing forces that factored character to be zero.

For a finite characteristic $\ell$-group quotient
$Q$ of $\widehat F$, the map to $\operatorname{Aut}(Q)$
factors through the actual quotient $P_\ell\to Q$.
Thus its quadratic characters pull back through
(B.F18). Finally $G_{\mathbb Q}\subseteq\GT$,
so every character in $\mathscr X$ has the required
arithmetic vanishing. This proves every stated scope.
\end{proof}

